\documentclass[spanish,letterpaper,oneside]{article}

\def\documenttitle {Aprendizaje Profundo y Cultura Escolar en Matemáticas}
\def\documentsubtitle {Fundamentos para la enseñanza y el aprendizaje I · Cohorte SD26}
\def\documentsubject {Actividad 5 · Bloque III (Módulo 51707)}
\def\documentauthor {Mtro. Martín Jonathan de la Cruz Muñoz}
\def\coursename {Fundamentos para la enseñanza y el aprendizaje I}
\def\coursecode {MAEM-SD26}

\def\universityname {Instituto de Investigación, Innovación y Estudios de Posgrado para la Educación del Estado de Nuevo León}
\def\universityfaculty {Maestría en Aprendizaje y Enseñanza de las Matemáticas}
\def\universitydepartment {Fundamentos para la enseñanza y el aprendizaje I}
\def\universitydepartmentimage {img/iiiepe/marca-iiiepe.png}
\def\universitydepartmentimagecfg {height=1.75cm}
\def\universitylocation {Monterrey, Nuevo León}
\def\studentfullname {Mtro. Martín Jonathan de la Cruz Muñoz}
\def\studentid {00874}
\def\studentemail {martin.delacruz@campus.iiiepe.edu.mx}
\def\programperiod {2do periodo, septiembre--diciembre de 2026 (SD26)}
\def\programmodality {Modalidad presencial}
\def\programshort {MAEM}

\def\authortable {
  \begin{tabular}{ll}
    \textbf{Alumno:} & \studentfullname \\
    Matrícula: & \studentid \\
    Correo institucional: & \studentemail \\
    Unidad: & \coursename \\
    Programa: & \programshort \\
    Periodo: & \programperiod \\
    & \\
    \multicolumn{2}{l}{Fecha de realización: 19 de noviembre de 2026} \\
    \multicolumn{2}{l}{\universitylocation}
  \end{tabular}
}

\input{template}

\usepackage{helvet}
\renewcommand{\familydefault}{\sfdefault}
\usepackage{array}
\usepackage{longtable}
\usepackage{pdflscape}
\usepackage{tikz}
\usepackage{pgfplots}
\pgfplotsset{compat=1.18}
\usepackage[most]{tcolorbox}
\usetikzlibrary{arrows.meta,positioning,calc,fit,backgrounds,shapes.geometric}

\definecolor{iiiepeNavy}{HTML}{163B5C}
\definecolor{iiiepeBlue}{HTML}{1F6F9F}
\definecolor{iiiepeAqua}{HTML}{20B7A8}
\definecolor{iiiepeOrange}{HTML}{F28C28}
\definecolor{iiiepePaper}{HTML}{F6F8F7}
\definecolor{iiiepeInk}{HTML}{1E2A32}
\definecolor{iiiepePale}{HTML}{EEF5F8}
\definecolor{iiiepeMint}{HTML}{E7F6F4}
\definecolor{iiiepeWarm}{HTML}{FFF3E6}
\def\portraittitlecolor {iiiepeNavy}
\setcitestyle{authoryear,open={(},close={)}}
\makeatletter
\renewcommand{\@biblabel}[1]{}
\makeatother
\setlength{\bibhang}{1.27cm}
\setlength{\bibsep}{6pt}
\AtBeginDocument{\renewcommand{\refname}{Referencias}}

\newcommand{\insertcoverwatermark}{%
  \AddToShipoutPictureBG*{%
    \begin{tikzpicture}[remember picture,overlay]
      \node[opacity=0.04,inner sep=0pt] at (current page.center) {%
        \includegraphics[width=0.65\paperwidth]{img/iiiepe/marca-iiiepe.png}%
      };
    \end{tikzpicture}%
  }%
}

\setlength{\parskip}{5pt}\setlength{\parindent}{0pt}
\begin{document}\insertcoverwatermark\templatePortrait\templatePagecfg\templateFinalcfg\fancyfoot{}\fancyfoot[L]{\scriptsize IIIEPE · MAEM · SD26}\fancyfoot[C]{\scriptsize\thepage}\fancyfoot[R]{\scriptsize Actividad 5 · Aprendizaje Profundo}\setstretch{1.3}\raggedbottom
\section{Introducción y postura central}
El tránsito de una enseñanza matemática centrada en repetir algoritmos hacia comunidades escolares capaces de comprender, argumentar y transferir conocimientos constituye uno de los desafíos neurálgicos de la educación contemporánea. A partir del análisis crítico de \citet{quiroga2022aprendizaje} sobre el aprendizaje profundo como herramienta para transformar la cultura escolar, articulado con los planteamientos de las ciencias del aprendizaje de \citet{sawyer2008optimizing} y \citet{fullan2014rich}, sostengo la tesis de que el dominio técnico de procedimientos matemáticos resulta insuficiente si no integra simultáneamente las dimensiones cognitiva, interpersonal e intrapersonal dentro de nuevas relaciones pedagógicas en el aula.

Mi postura es que las dificultades habituales observadas en bachillerato técnico e ingeniería —cuando el alumnado ejecuta reglas algebraicas mecánicas sin explicar su sentido físico ni adaptar sus estrategias ante situaciones nuevas— revelan el predominio de una cultura escolar orientada hacia el aprendizaje superficial. Por ello, este informe responde de manera fundamentada a las cinco preguntas solicitadas en el módulo \texttt{51707}, integrando un diagrama triádico en TikZ, un contraste formal entre destreza procedimental y comprensión profunda mediante un ejemplo matemático concreto, y orientaciones de transferencia docente en coherencia con las nuevas pedagogías.

\section{Desarrollo argumentado de las cinco interrogantes centrales}

\subsection*{1. ¿Qué diferencias fundamentales existen entre el aprendizaje superficial y el aprendizaje profundo?}
Retomando la distinción clásica de \citet{marton1976qualitative} analizada por \citet{quiroga2022aprendizaje}, el aprendizaje superficial consiste en abordar las tareas escolares con la intención mínima de cumplir requisitos externos mediante memorización mecánica, reproducción literal de definiciones y tratamiento aislado de contenidos sin buscar conexiones globales. En contraste, el aprendizaje profundo persigue comprender relaciones significativas, articular conceptos nuevos con saberes previos, examinar evidencias de manera crítica y transferir lo comprendido a problemas inéditos \citep{sawyer2008optimizing}.

\begin{tcolorbox}[colback=iiiepePale,colframe=iiiepeNavy,title=\textbf{Síntesis comparativa · Enfoque superficial frente a aprendizaje profundo},breakable]
{\small\renewcommand{\arraystretch}{1.28}
\begin{tabular}{|p{0.23\linewidth}|p{0.35\linewidth}|p{0.35\linewidth}|}\hline
\rowcolor{iiiepePale} \textbf{Criterio de análisis} & \textbf{Aprendizaje superficial} & \textbf{Aprendizaje profundo} \\\hline
Intención estudiantil & Cumplir la tarea o aprobar la evaluación con mínimo esfuerzo cognitivo \citep{marton1976qualitative}. & Comprender causas, conexiones estructurales y significado aplicado \citep{quiroga2022aprendizaje}. \\\hline
Tratamiento del saber & Fragmentado, memorístico y dependiente de señales literales del enunciado. & Relacional, crítico y transferible entre representaciones múltiples \citep{sawyer2008optimizing}. \\\hline
Uso del error & Se oculta o penaliza sin análisis conceptual posterior. & Se analiza colectivamente para revisar supuestos y reconstruir argumentos \citep{eef2021feedback}. \\\hline
\end{tabular}}
\end{tcolorbox}

\subsection*{2. ¿Qué características de la enseñanza tradicional de las matemáticas favorecen un aprendizaje superficial?}
El análisis propio de la enseñanza tradicional en matemáticas muestra cuatro rasgos que inducen pasividad y mecanización: (a) presentar reglas y algoritmos acabados antes de permitir explorar el fenómeno o problema que les dio origen; (b) proponer listas extensas de ejercicios idénticos que premian la velocidad repetitiva por encima del razonamiento; (c) fragmentar los contenidos —tablas, gráficas, álgebra y contexto físico— como capítulos inconexos; y (d) concentrar la evaluación en respuestas numéricas únicas sin requerir justificación verbal o escrita de las decisiones tomadas \citep{star2015algebra,eef2017mathematics}.

\subsection*{3. ¿Por qué el dominio de procedimientos y algoritmos matemáticos no garantiza aprendizaje profundo?}
Ejecutar con destreza un algoritmo garantiza únicamente fluidez operatoria dentro de formatos rutinarios conocidos; sin embargo, no asegura comprender qué modelan las variables, bajo qué condiciones aplica el procedimiento ni cómo validar si la respuesta tiene sentido real \citep{sawyer2008optimizing}. Considérese un ejemplo representativo en bachillerato técnico e ingeniería: un estudiante puede resolver mecánicamente la ecuación lineal
\[
0.5x + 2 = 0.2x + 3.8 \implies 0.3x = 1.8 \implies x = 6
\]
aplicando correctamente reglas de transposición de términos. No obstante, cuando esa misma relación representa el punto en que dos sensores térmicos o dos planes de mantenimiento igualan su costo o lectura ($y_1 = 0.5x + 2$ frente a $y_2 = 0.2x + 3.8$), el aprendizaje superficial se evidencia si no distingue que $2$ y $3.8$ son condiciones iniciales (sesgos u ordenadas al origen), que ambas relaciones son afines y no proporcionales ($y=kx$), ni interpreta en qué intervalo conviene un esquema sobre otro \citep{conalep2025mecatronica}. Esto implica que la fluidez algorítmica es condición útil pero nunca suficiente para la comprensión matemática profunda.

\subsection*{4. ¿Cómo se relacionan las dimensiones cognitiva, interpersonal e intrapersonal con las matemáticas?}
Como sistematizan \citet{quiroga2022aprendizaje}, el aprendizaje profundo exige la convergencia inseparable de tres dimensiones formativas (Figura~\ref{fig:act5-triada}). En matemáticas, la dimensión \textbf{cognitiva} abarca la comprensión relacional, el pensamiento lógico, el uso estratégico de representaciones y la resolución de problemas no rutinarios; la dimensión \textbf{interpersonal} se activa al argumentar conjeturas frente a pares, escuchar estrategias alternativas y construir acuerdos verificables; finalmente, la dimensión \textbf{intrapersonal} involucra la autorregulación ante el bloqueo inicial, la autoeficacia y la mentalidad de crecimiento para asumir el error como insumo de mejora y no como incapacidad fija \citep{schunk2012teorias}.

\begin{figure}[htbp]
\centering
\begin{tikzpicture}[font=\scriptsize, >=Latex]
  \node[draw=iiiepeNavy,fill=iiiepeNavy!12,very thick,circle,minimum size=2.45cm,text width=2.2cm,align=center,font=\scriptsize\bfseries,text=iiiepeNavy] (cen) at (0,0) {APRENDIZAJE\\PROFUNDO EN\\MATEMÁTICAS};

  \node[draw=iiiepeNavy,fill=iiiepePale,thick,rounded corners=4pt,text width=5.6cm,align=left,inner sep=6pt] (cog) at (0,3.2) {%
    \textbf{\color{iiiepeNavy}1. Dimensión Cognitiva}\\
    Comprensión conceptual, pensamiento crítico, modelación algebraica/gráfica, análisis de límites de validez y transferencia a problemas auténticos \citep{sawyer2008optimizing}.
  };

  \node[draw=iiiepeBlue,fill=iiiepePale,thick,rounded corners=4pt,text width=5.2cm,align=left,inner sep=6pt] (inter) at (-4.8,-2.4) {%
    \textbf{\color{iiiepeBlue}2. Dimensión Interpersonal}\\
    Colaboración deliberada en equipos, comunicación matemática clara, defensa argumentada de hipótesis y escucha respetuosa de rutas distintas \citep{quiroga2022aprendizaje}.
  };

  \node[draw=iiiepeAqua!80!black,fill=iiiepeMint,thick,rounded corners=4pt,text width=5.2cm,align=left,inner sep=6pt] (intra) at (4.8,-2.4) {%
    \textbf{\color{iiiepeAqua!75!black}3. Dimensión Intrapersonal}\\
    Autorregulación metacognitiva, perseverancia ante tareas complejas, autoeficacia (Bandura) y mentalidad de crecimiento frente al error \citep{schunk2012teorias}.
  };

  \draw[<->,very thick,iiiepeNavy] (cog.south west) -- (inter.north east);
  \draw[<->,very thick,iiiepeAqua!80!black] (cog.south east) -- (intra.north west);
  \draw[<->,very thick,iiiepeOrange!85!black] (inter.east) -- (intra.west);
  \draw[->,dashed,thick,iiiepeNavy] (cog.south) -- (cen.north);
  \draw[->,dashed,thick,iiiepeBlue] (inter.north east) -- (cen.south west);
  \draw[->,dashed,thick,iiiepeAqua!80!black] (intra.north west) -- (cen.south east);
\end{tikzpicture}
\caption{Interdependencia de las dimensiones cognitiva, interpersonal e intrapersonal en el aprendizaje matemático.}
\label{fig:act5-triada}
\end{figure}

\subsection*{5. ¿Qué transformaciones requiere la práctica docente para favorecer el aprendizaje profundo?}
De acuerdo con \citet{quiroga2022aprendizaje} y \citet{fullan2014rich}, transformar la docencia exige actuar simultáneamente en tres niveles de la cultura escolar: (1) los \textbf{recursos} de trabajo, sustituyendo guías mecánicas por retos abiertos vinculados con el entorno; (2) la \textbf{didáctica}, transitando desde el control unidireccional hacia relaciones de asociación entre docente y estudiantes que habiliten progresivamente la autonomía; y (3) las \textbf{creencias pedagógicas} compartidas sobre qué significa una buena clase de matemáticas. En paralelo, requiere instalar evaluación formativa oportuna centrada en criterios transparentes de comprensión y revisión reflexiva \citep{eef2021feedback,bostrom2023formative}.

{\small\setstretch{1.15}\renewcommand{\arraystretch}{1.32}
\begin{longtable}{|>{\raggedright\arraybackslash}p{0.22\linewidth}|>{\raggedright\arraybackslash}p{0.36\linewidth}|>{\raggedright\arraybackslash}p{0.36\linewidth}|}\hline
\rowcolor{iiiepePale} \textbf{Eje de cambio docente} & \textbf{Inercia tradicional} & \textbf{Transformación hacia las nuevas pedagogías} \\\hline
\endfirsthead
\hline
\rowcolor{iiiepePale} \textbf{Eje de cambio docente} & \textbf{Inercia tradicional} & \textbf{Transformación hacia las nuevas pedagogías} \\\hline
\endhead
Rol del docente & Expositor único de algoritmos y verificador final de respuestas. & Diseñador de desafíos, activador cognitivo y socio en la indagación \citep{fullan2014rich}. \\\hline
Rol del estudiante & Ejecutor individual de pasos repetidos bajo instrucción cerrada. & Co-constructor de explicaciones, interlocutor crítico y agente autónomo \citep{quiroga2022aprendizaje}. \\\hline
Evaluación en el aula & Calificación punitiva centrada en aciertos numéricos aislados. & Devolución formativa descriptiva sobre modelación, supuestos y errores \citep{bostrom2023formative}. \\\hline
Cultura escolar & Trabajo docente aislado y cobertura acelerada del programa. & Indagación colegiada sobre evidencia de aprendizaje y creencias compartidas \citep{quiroga2022aprendizaje}. \\\hline
\end{longtable}}

\section{Transferencia profesional al aula técnica e ingeniería}
En mi práctica cotidiana en bachillerato técnico e ingeniería, favorecer el aprendizaje profundo implica reorganizar el tiempo de clase para que el algoritmo deje de ser el punto de partida y pase a ser una herramienta justificada al servicio de problemas reales de mantenimiento, circuitos o análisis de datos. Además, integrar roles colaborativos explícitos y preguntas metacognitivas en cada entrega fortalece la autoeficacia del alumnado y reduce el abandono silencioso ante la complejidad matemática \citep{conalep2025mecatronica}.

\clearpage
\section{Conclusiones y síntesis argumentativa}
En conclusión, el análisis de \citet{quiroga2022aprendizaje} demuestra que abandonar el aprendizaje superficial en matemáticas no se resuelve multiplicando tareas repetitivas ni modernizando únicamente herramientas de presentación. Requiere articular rigor cognitivo, diálogo interpersonal y autorregulación intrapersonal bajo una cultura docente dispuesta a revisar sus propias creencias sobre enseñar y evaluar.

Por tanto, asumir el enfoque del aprendizaje profundo en esta cohorte SD26 compromete mi labor docente a diseñar secuencias basadas en asociación colaborativa, argumentación conceptual y evaluación formativa auténtica, sentando las bases directas para la propuesta de innovación didáctica que se diseña en la actividad siguiente.

\clearpage
\begin{thebibliography}{99}
\bibitem[Boström y Palm(2023)]{bostrom2023formative} Boström, E., y Palm, T. (2023). Formative assessment practice in mathematics: Teachers' perceptions of prerequisite conditions. \textit{Frontiers in Education}, 8, Artículo 1101192.
\bibitem[Collin y Quigley(2021)]{eef2021feedback} Collin, J., y Quigley, A. (2021). \textit{Teacher feedback to improve pupil learning}. Education Endowment Foundation.
\bibitem[CONALEP(2025)]{conalep2025mecatronica} Colegio Nacional de Educación Profesional Técnica. (2025). \textit{Perfil de egreso: Profesional Técnico-Bachiller en Mecatrónica. Modelo Académico 2025}. Secretaría Académica.
\bibitem[Fullan y Langworthy(2014)]{fullan2014rich} Fullan, M., y Langworthy, M. (2014). \textit{A rich seam: How new pedagogies find deep learning}. Pearson.
\bibitem[Henderson et~al.(2017)]{eef2017mathematics} Henderson, P., Hodgen, J., Foster, C., y Kuchemann, D. (2017). \textit{Improving mathematics in Key Stages 2 and 3}. Education Endowment Foundation.
\bibitem[Marton y Säljö(1976)]{marton1976qualitative} Marton, F., y Säljö, R. (1976). On qualitative differences in learning: I. Outcome and process. \textit{British Journal of Educational Psychology}, 46(1), 4--11.
\bibitem[Quiroga Aguilar y Lara Órdenes(2022)]{quiroga2022aprendizaje} Quiroga Aguilar, L., y Lara Órdenes, E. (2022). El aprendizaje profundo como herramienta para cambio en la visión de aprendizaje de una cultura escolar. \textit{Revista Educación las Américas}, 12(1). \url{https://doi.org/10.35811/rea.v12i1.201}
\bibitem[Sawyer(2008)]{sawyer2008optimizing} Sawyer, R. K. (2008). Optimizing learning: Implications of learning sciences research. OECD Publishing.
\bibitem[Schunk(2012)]{schunk2012teorias} Schunk, D. H. (2012). \textit{Teorías del aprendizaje: Una perspectiva educativa} (6.ª ed.). Pearson Educación.
\bibitem[Star et~al.(2015)]{star2015algebra} Star, J. R., Caronongan, P., Foegen, A., Furgeson, J., Keating, B., Larson, M. R., Lyskawa, J., McCallum, W. G., Porath, J., y Zbiek, R. M. (2015). \textit{Teaching strategies for improving algebra knowledge}. U.S. Department of Education.
\end{thebibliography}

\end{document}
